% Part:sets-functions-relations
% Chapter: size-of-sets
% Section: enumerations

\documentclass[../../../include/open-logic-section]{subfiles}

\begin{document}

\olfileid{sfr}{siz}{enm}

\olsection{प्रगणने आणि \usetoken{S}{enumerable} संच}

\begin{editorial}
  या विभागात संचांच्या प्रगणनांची चर्चा आहे. त्यांची व्याख्या
  $\PosInt$ पासूनच्या आच्छादनांनी केली आहे. गणिताची फारशी पार्श्वभूमी
  नसलेल्या वाचकांसाठी विषय हळूहळू मांडला आहे. पर्यायी, अधिक संक्षिप्त
  आवृत्ती \olref[enm-alt]{sec} मध्ये दिली आहे. तेथे प्रगणनांची वेगळी
  व्याख्या केली आहे: $\Nat$ किंवा त्याच्या आरंभीच्या खंडाशी
  एकास-एक आच्छादन असणे.
\end{editorial}

\begin{explain}
संचांचे !!{element}s यादीत देऊन आपण आधीच संचांची उदाहरणे दिली आहेत.
संच अनंत असला, तरी त्याच्या !!{element}s ची यादी कशी आणि कधी
देता येते, याची आता अधिक व्यापकपणे चर्चा करू.
\end{explain}

\begin{defn}[प्रगणनाची अनौपचारिक व्याख्या]
अनौपचारिकपणे, संच~$A$ चे \emph{प्रगणन} म्हणजे~$A$ मधील
!!{element}s ची अशी यादी (कदाचित अनंत) की $A$ मधील प्रत्येक
!!{element} त्या यादीत कोणत्या तरी सांत क्रमांकाच्या स्थानी येतो.
$A$ चे प्रगणन असेल, तर $A$ हा \emph{!!{enumerable}} आहे असे म्हणतात.
\end{defn}

\begin{explain}
प्रगणनांविषयी काही मुद्दे:
\begin{enumerate}
\item ज्या यादीला सुरुवात असते आणि जिच्यात पहिल्या घटकाव्यतिरिक्त
  प्रत्येक !!{element} च्या लगेच आधी एकच !!{element} असतो, अशीच यादी
  आपण प्रगणन मानतो. दुसऱ्या शब्दांत, यादीच्या पहिल्या !!{element}
  आणि इतर कोणत्याही !!{element} यांच्यामध्ये केवळ सांत संख्येने घटक
  असतात. विशेषतः, प्रगणनातील प्रत्येक !!{element} चे स्थान सांत
  क्रमांकाचे असते: पहिल्या !!{element} चे स्थान~$1$, दुसऱ्याचे~$2$, इत्यादी.
\item एकाच संच~$A$ ची वेगवेगळी प्रगणने असू शकतात, ज्यांत
  !!{element}s येण्याचा क्रम वेगळा असतो: $4$, $1$, $25$, $16$,~$9$
  हे पहिल्या पाच वर्गसंख्यांच्या संचाचे प्रगणन आहे, तसेच
  $1$, $4$, $9$, $16$,~$25$ हेही आहे.
\item पुनरुक्ती असलेली प्रगणनेही प्रगणनेच असतात: $1$, $1$, $2$,
  $2$, $3$, $3$,~\dots{} हे त्याच संचाचे प्रगणन आहे, ज्याचे
  प्रगणन $1$, $2$, $3$,~\dots{} हे आहे.
\item विशिष्ट प्रगणन सांगताना क्रम आणि पुनरुक्ती यांना महत्त्व
  \emph{असते}: धन पूर्णांकांचे प्रगणन $1$, $2$, $3$, $1$, \dots{}
  अशा सुरुवातीने करता येते; पण नेहमीच्या $1$, $2$, $3$, $4$,~\dots
  या पद्धतीने प्रगणन केले, तर त्यातील नमुना अधिक सहज दिसतो.
\item प्रगणनाला सुरुवात असली पाहिजे: \dots, $3$, $2$, $1$ हे
  धन पूर्णांकांचे प्रगणन नाही, कारण त्याला पहिला !!{element} नाही.
  अनौपचारिक व्याख्येवरून हे कसे येते ते पाहण्यासाठी स्वतःला विचारा:
  ``या यादीत 76 ही संख्या कोणत्या स्थानी येते?''
\item पुढील यादी धन पूर्णांकांचे प्रगणन नाही:
  $1$, $3$, $5$, \dots, $2$, $4$, $6$, \dots\@ अडचण अशी की
  सम संख्या सांत क्रमांकांच्या स्थानांवर येण्याऐवजी
  $\infty + 1$, $\infty + 2$, $\infty + 3$ या स्थानांवर येतात.
\item रिक्त संच गणनीय आहे: रिक्त यादी हे त्याचे प्रगणन आहे!{}
\end{enumerate}
\end{explain}

\begin{prop}
  $A$ चे प्रगणन असेल, तर त्याचे पुनरुक्ती नसलेले प्रगणनही असते.
\end{prop}

\begin{proof}
  $A$ चे $x_1$, $x_2$, \dots{} असे प्रगणन आहे, आणि प्रत्येक
  $x_i$ हा~$A$ चा !!{element} आहे असे समजा. प्रगणनातून पुन्हा
  आलेले !!{element}s काढून टाकून त्यातील पुनरुक्ती काढता येते.
  उदाहरणार्थ, पुढीलप्रमाणे नवे प्रगणन मिळवता येते: $x_i$ हा~$A$
  मधील !!a{element} असून $x_1$, \dots, $x_{i-1}$ यांमध्ये नसेल,
  तर त्याला यादीत लिहायचे; आणि $x_i$ हा आधीच $x_1$, \dots,~$x_{i-1}$
  यांमध्ये असेल, तर त्याला यादीतून काढून टाकायचे.
\end{proof}

मागील युक्तिवाद दाखवतो की प्रगणने आणि !!{enumerable} संच नीट
समजून घेण्यासाठी, तसेच त्यांविषयी गोष्टी सिद्ध करण्यासाठी,
अधिक नेमकी व्याख्या हवी. ती पुढे दिली आहे.

\begin{defn}[प्रगणनाची औपचारिक व्याख्या]
$A \neq \emptyset$ या संचाचे \emph{प्रगणन} म्हणजे कोणतेही
!!{surjective} फलन $f \colon \PosInt \to A$.
\end{defn}

\begin{explain}
औपचारिक व्याख्या आणि कदाचित अनंत असणाऱ्या यादीचा वापर करणारी
अनौपचारिक व्याख्या समतुल्य आहेत याची खात्री करू. प्रथम, $\PosInt$
पासून संच~$A$ पर्यंतचे कोणतेही !!{surjective} फलन~$A$ चे प्रगणन करते.
अशा फलनाने वरील अनौपचारिक अर्थाचे प्रगणन ठरते: $f(1)$, $f(2)$,
$f(3)$, \dots{} ही यादी. $f$ हे !!{surjective} असल्यामुळे~$A$
मधील प्रत्येक !!{element} हा कोणत्या तरी~$n \in \PosInt$ साठीचे
~$f(n)$ चे मूल्य असतो. म्हणून $A$ मधील प्रत्येक !!{element} यादीत
सांत क्रमांकाच्या स्थानी येतो. फलन !!{injective} असेलच असे नसल्याने
यादीत पुनरुक्ती असू शकते; पण वर म्हटल्याप्रमाणे ती चालते.

उलट,~$A$ मधील सर्व !!{element}s चे प्रगणन करणारी यादी दिली असल्यास
!!a{surjective} फलन $f\colon \PosInt \to A$ असे परिभाषित करता येते:
$f(n)$ म्हणजे यादीतील $n$ व्या स्थानी असलेला !!{element}; आणि
$n$ व्या स्थानी !!{element} नसेल, तर यादीचा शेवटचा !!{element}.
$A$ रिक्त असतो आणि म्हणून यादीही रिक्त असते, फक्त या एकाच बाबतीत
या पद्धतीने !!a{surjective} फलन मिळत नाही. म्हणून प्रत्येक अरिक्त
यादी एक !!a{surjective} फलन $f\colon \PosInt \to A$ ठरवते.
\end{explain}

\begin{defn}
  \ollabel{defn:enumerable}
  संच~$A$ हा !!{enumerable} असतो तेव्हा आणि केवळ तेव्हाच,
  जेव्हा तो रिक्त असतो किंवा त्याचे प्रगणन असते.
\end{defn}

\begin{ex}
धन पूर्णांकांचे ($\PosInt$) प्रगणन करणारे फलन म्हणजे फक्त
$f(n) = n$ ने दिलेले एकरूपता फलन. नैसर्गिक संख्या $\Nat$ यांचे
प्रगणन करणारे फलन म्हणजे $g(n) = n - 1$.
\end{ex}

\begin{ex}
$f\colon \PosInt \to \PosInt$ आणि $g \colon \PosInt \to \PosInt$
ही फलने अशी दिली आहेत:
\begin{align*}
f(n) & = 2n \text{ आणि}\\
g(n) & = 2n - 1
\end{align*}
ती अनुक्रमे सम धन पूर्णांकांचे आणि विषम धन पूर्णांकांचे प्रगणन करतात.
मात्र यांपैकी कोणतेही फलन $\PosInt$ चे प्रगणन नाही, कारण कोणतेही
!!{surjective} नाही.
\end{ex}

\begin{prob}
धन वर्गसंख्या $1$, $4$, $9$, $16$, \dots{} यांचे प्रगणन परिभाषित करा.
\end{prob}

\begin{ex}
$f(n) = (-1)^{n} \lceil \frac{(n-1)}{2}\rceil$ हे फलन पूर्णांकांचा
संच~$\Int$ प्रगणित करते. येथे $\lceil x \rceil$ हे
\emph{ऊर्ध्व पूर्णांक फलन} दर्शवते; ते $x$ ला त्याच्यापेक्षा कमी
नसलेल्या सर्वांत जवळच्या पूर्णांकाकडे नेते. धन आणि ऋण पूर्णांकांमध्ये
आलटूनपालटून ``उड्या मारून'' $f$ हे $\Int$ ची मूल्ये कशी निर्माण
करते ते पाहा:
\[
\begin{array}{c c c c c c c c}
f(1) & f(2) & f(3) & f(4) & f(5) & f(6) & f(7) & \dots \\ \\
- \lceil \tfrac{0}{2} \rceil & \lceil \tfrac{1}{2}\rceil & - \lceil \tfrac{2}{2} \rceil & \lceil \tfrac{3}{2} \rceil & - \lceil \tfrac{4}{2} \rceil  & \lceil \tfrac{5}{2}
\rceil & - \lceil \tfrac{6}{2} \rceil & \dots \\ \\
0 & 1 & -1 & 2 & -2 & 3 & -3 & \dots
\end{array}
\]
%READERNOTE{OLSIZ-001: गोठवलेल्या इंग्रजी तक्त्याच्या शेवटच्या ओळीत f(7) साठी अपेक्षित −3 हे मूल्य सुटले आहे; सूत्र आणि वरील दोन ओळी ते मूल्य निश्चित करतात. मराठी तक्त्यात −3 समाविष्ट करून ही उणीव दुरुस्त केली आहे.}
$f$ ची व्याख्या पुढीलप्रमाणे प्रकरणे पाडून केली आहे असेही मानता येते:
\[
f(n) = \begin{cases}
  0 & \text{जर $n = 1$ असेल}\\
  n/2 & \text{जर $n$ सम असेल}\\
  -(n-1)/2 & \text{जर $n$ विषम आणि $>1$ असेल}
  \end{cases}
\]
\end{ex}

\begin{prob}
  $A$ आणि $B$ हे !!{enumerable} असतील, तर $A \cup B$ हाही
  गणनीय असतो, हे दाखवा. यासाठी !!{surjective} फलने
  $f\colon \PosInt \to A$ आणि $g\colon \PosInt \to B$ आहेत असे
  समजा. एक !!a{surjective} फलन~$h\colon \PosInt \to A \cup B$
  परिभाषित करा आणि ते !!{surjective} आहे हे सिद्ध करा. $A$ किंवा
  ~$B = \emptyset$ असण्याच्याही बाबतींचा विचार करा.
\end{prob}

\begin{prob}
  $B \subseteq A$ आणि $A$ हा !!{enumerable} असेल, तर~$B$ हाही
  गणनीय असतो, हे दाखवा. यासाठी एक !!a{surjective} फलन
  $f\colon \PosInt \to A$ आहे असे समजा. एक !!a{surjective}
  फलन~$g\colon \PosInt \to B$ परिभाषित करा आणि ते
  !!{surjective} आहे हे सिद्ध करा. $B = \emptyset$ असल्यास काय होते?
\end{prob}

\begin{prob}
  $A_1$, $A_2$, \dots, $A_n$ हे सर्व !!{enumerable} असतील, तर
  $A_1 \cup \dots \cup A_n$ हाही गणनीय असतो, हे
  $n$ वर विगमन करून दाखवा. $A$ आणि~$B$ हे दोन संच
  !!{enumerable} असल्यास~$A \cup B$ हाही गणनीय असतो
  हे तथ्य तुम्ही गृहीत धरू शकता.
\end{prob}

संचातील !!{element}s ची यादी करताना $1$ व्या !!{element} पासून
मोजायला सुरुवात करणे कदाचित अधिक स्वाभाविक असले, तरी गणितज्ञांना
मोजणीसाठी नैसर्गिक संख्या~$\Nat$ वापरायला आवडतात. ते यादीतील
$0$ व्या, $1$ व्या, $2$ व्या, अशा !!{element}s विषयी बोलतात.
त्याला अनुसरून प्रगणनाची व्याख्या $\Nat$ पासून~$A$ पर्यंतचे
!!a{surjective} फलन अशी करता येते. अर्थात, या दोन व्याख्या समतुल्य आहेत.

\begin{prop}\ollabel{prop:enum-shift}
  !!a{surjection} $f\colon \PosInt \to A$ असते तेव्हा आणि
  केवळ तेव्हाच, जेव्हा !!a{surjection} $g\colon \Nat \to A$ असते.
\end{prop}

\begin{proof}
  !!a{surjection} $f\colon \PosInt \to A$ दिले असता, सर्व
  $n \in \Nat$ साठी $g(n) = f(n+1)$ अशी व्याख्या करता येते.
  $g\colon \Nat \to A$ हे !!{surjective} आहे हे सहज दिसते.
  उलट, !!a{surjection} $g\colon \Nat \to A$ दिले असता,
  $f(n) = g(n-1)$ अशी व्याख्या करा.
\end{proof}

यावरून पुढील निष्कर्ष मिळतो:

\begin{cor}\ollabel{cor:enum-nat}
संच $A$ हा !!{enumerable} असतो तेव्हा आणि केवळ तेव्हाच, जेव्हा
तो रिक्त असतो किंवा !!a{surjective} फलन $f\colon \Nat \to A$ असते.
\end{cor}

संच~$A$ मधील !!{element}s च्या यादीतून पुनरुक्ती नसलेली यादी
मिळवता येते याची वर चर्चा केली. प्रगणनांबाबतही हे खरे आहे,
पण ते काटेकोरपणे मांडणे आणि सिद्ध करणे थोडे कठीण आहे. कोणतेही
फलन $f\colon \PosInt \to A$ हे सर्व $n \in \PosInt$ साठी
परिभाषित असले पाहिजे. ~$A$ मध्ये फक्त सांत संख्येने !!{element}s
असतील, तर~$\PosInt$ मधील अनंत संख्येने असलेल्या !!{element}s
वर परिभाषित असणारे आणि~$A$ मधील सर्व !!{element}s मूल्ये म्हणून
घेणारे, पण एकच मूल्य दोनदा कधीच न घेणारे फलन असू शकत नाही.
अशा बाबतीत, म्हणजे पुनरुक्ती नसलेली यादी सांत असते तेव्हा,~$f$
साठी वेगळा प्रांत निवडला पाहिजे; त्यात केवळ सांत संख्येने
!!{element}s असले पाहिजेत. पुनरुक्ती नसणे म्हणजे $f$ हे
!!{injective} असले पाहिजे. ते !!{surjective} ही असल्यामुळे
आपण कोणता तरी सांत संच $\{1, \dots, n\}$ किंवा $\PosInt$
आणि~$A$ यांच्यामधील !!a{bijection} शोधत आहोत.

\begin{prop}\ollabel{prop:enum-bij}
$f\colon \PosInt \to A$ हे !!{surjective} असेल (म्हणजे~$A$
चे प्रगणन असेल), तर !!a{bijection} $g\colon Z \to A$ असते,
जेथे $Z$ हा~$\PosInt$ किंवा कोणत्या तरी~$n \in \PosInt$
साठी $\{1, \dots, n\}$ असतो.
\end{prop}

\begin{proof}
  फलन $g$ ची व्याख्या आपण पुनरावर्ती रीतीने करतो: $g(1) = f(1)$
  असे ठेवा. $g(i)$ आधीच परिभाषित झाले असेल, तर $f(1)$, $f(2)$,
  \dots{} यांपैकी $g(1)$, \dots, $g(i)$ यांत आधीच न आलेले
  पहिले मूल्य, असे मूल्य असेल तर, $g(i+1)$ म्हणून घ्या.
  $A$ मध्ये नेमके $n$ !!{element}s असतील, तर $g(1)$, \dots, $g(n)$
  ही सर्व मूल्ये परिभाषित असतात, आणि आपल्याला
  $g\colon \{1, \dots, n\} \to A$ असे फलन मिळते.
  $A$ मध्ये अनंत संख्येने !!{element}s असतील, तर कोणत्याही $i$
  साठी $f(1)$, $f(2)$, \dots{} या प्रगणनात~$A$ मधील असा
  !!a{element} असलाच पाहिजे, जो $g(1)$, \dots, $g(i)$ यांत
  आधीच आलेला नाही. अशा बाबतीत आपण $g\colon \PosInt \to A$
  हे फलन परिभाषित केले आहे.

  फलन $g$ हे !!{surjective} आहे, कारण~$A$ मधील कोणताही घटक
  $f(1)$, $f(2)$, \dots{} यांमध्ये येतो ($f$ हे !!{surjective}
  असल्यामुळे), आणि म्हणून कधी ना कधी तो कोणत्या तरी~$i$
  साठीचे~$g(i)$ चे मूल्य होईल. ते !!{injective} ही आहे:
  $g(j) = g(i)$ असणारे $j < i$ असते, तर $g(i)$ हे मूल्य
  आधीच $g(1)$, \dots, $g(i-1)$ यांमध्ये आलेले असते;
  पण हे~$g$ च्या आपल्या व्याख्येविरुद्ध आहे.
\end{proof}

\begin{cor}\ollabel{cor:enum-nat-bij}
संच $A$ हा !!{enumerable} असतो तेव्हा आणि केवळ तेव्हाच, जेव्हा
तो रिक्त असतो किंवा !!a{bijection} $f\colon N \to A$ असते,
जेथे $N = \Nat$ किंवा कोणत्या तरी $n \in \Nat$ साठी
$N = \{0, \dots, n\}$ असते.
\end{cor}

\begin{proof}
$A$ हा !!{enumerable} असतो तेव्हा आणि केवळ तेव्हाच, जेव्हा $A$
रिक्त असतो किंवा !!a{surjective} $f\colon \PosInt \to A$ असते.
\olref{prop:enum-bij} नुसार दुसरी अट तेव्हा आणि केवळ तेव्हाच
खरी असते, जेव्हा !!a{bijective} फलन~$f\colon Z \to A$ असते,
जेथे $Z = \PosInt$ किंवा कोणत्या तरी $n \in \PosInt$ साठी
$Z = \{1, \dots, n\}$ असते. \olref{prop:enum-shift} च्या
सिद्धतेतील युक्तिवादानुसार, हे तेव्हा आणि केवळ तेव्हाच शक्य होते,
जेव्हा !!a{bijection} $g\colon N \to A$ असते, जेथे
$N = \Nat$ किंवा $N = \{0, \dots, n-1\}$ असते.
%READERNOTE{MRSIZ-001: विधानातील सांत प्रांत 0 ते n आणि सिद्धतेतील 0 ते n−1 हे वेगवेगळ्या चलनामांनी त्याच सांत आरंभीच्या खंडांच्या कुटुंबाचे वर्णन करतात; येथे स्रोताचे दोन्ही निर्देशांक जसेच्या तसे ठेवले आहेत.}
\end{proof}

\begin{prob}
  \olref[sfr][siz][enm]{defn:enumerable} नुसार संच $A$ हा गणनीय
  असतो तेव्हा आणि केवळ तेव्हाच, जेव्हा $A = \emptyset$ किंवा
  !!a{surjective} $f\colon \PosInt \to A$ असते.
  ``!!{enumerable} संच'' याची नेमकी व्याख्या पुढीलप्रमाणेही करता येते:
  संच गणनीय असतो तेव्हा आणि केवळ तेव्हाच, जेव्हा !!a{injective}
  फलन $g\colon A \to \PosInt$ असते. या व्याख्या समतुल्य आहेत
  हे दाखवा. म्हणजे, !!a{injective} फलन $g\colon A \to \PosInt$
  असते तेव्हा आणि केवळ तेव्हाच, जेव्हा $A = \emptyset$ किंवा
  !!a{surjective} $f\colon \PosInt \to A$ असते, हे दाखवा.
\end{prob}

\end{document}
